\chapter{Proof Techniques}
\signature

\section{The vocabulary of mathematical statements}

A \emph{theorem} is a statement proved true; a \emph{lemma} is a helper
theorem; a \emph{corollary} follows easily from a theorem; a
\emph{conjecture} awaits proof; a \emph{proof} is a valid argument, built
from axioms, definitions and previous results by the rules of inference of
Chapter~5. Most theorems have the form
$\forall x\,\bigl(P(x) \to Q(x)\bigr)$, and the techniques below are
strategies for establishing $P \to Q$ for an arbitrary $x$ (whence UG
concludes).

\section{Direct proof}

To prove $p \to q$ directly: assume $p$, derive $q$ by definitions and known
results.

\begin{theorem}
If $n$ is odd then $n^{2}$ is odd.
\end{theorem}

\begin{proof}
Let $n = 2k + 1$ for some integer $k$. Then
$n^{2} = 4k^{2} + 4k + 1 = 2(2k^{2} + 2k) + 1$, which is odd.
\end{proof}

\section{Proof by contraposition}

To prove $p \to q$, prove $\neg q \to \neg p$ instead --- the two are
equivalent (Chapter~2).

\begin{theorem}
If $n^{2}$ is even then $n$ is even.
\end{theorem}

\begin{proof}
Contrapositive: if $n$ is odd then $n^{2}$ is odd --- proved above.
\end{proof}

Choose contraposition when $\neg q$ is more concrete to work with than $p$
(here: ``$n$ odd'' gives an explicit form, ``$n^2$ even'' does not).

\section{Proof by contradiction}

To prove $\varphi$: assume $\neg\varphi$ and derive a contradiction
($r \wedge \neg r$); then $\neg\varphi$ is impossible, so $\varphi$ holds.
\emph{Template:} ``Suppose, for contradiction, that \ldots\ Then \ldots\
contradiction. Hence \ldots''

\begin{theorem}[Canonical example]\label{thm:sqrt2}
$\sqrt{2}$ is irrational.
\end{theorem}

\begin{proof}
Suppose $\sqrt2 = a/b$ with $a, b \in \Z^{+}$ in lowest terms. Then
$a^{2} = 2b^{2}$, so $a^{2}$ is even, so $a$ is even (previous theorem):
$a = 2c$. Then $4c^{2} = 2b^{2}$, i.e.\ $b^{2} = 2c^{2}$, so $b$ is even
too --- contradicting lowest terms. Hence $\sqrt2$ is irrational.
\end{proof}

\begin{theorem}[Euclid]\label{thm:euclid}
There are infinitely many primes.
\end{theorem}

\begin{proof}
Suppose $p_1, \ldots, p_k$ were all the primes and set
$N = p_1 \cdots p_k + 1$. Each $p_i$ divides $N - 1$, so none divides $N$;
yet $N > 1$ has a prime divisor --- a prime outside the list:
contradiction.
\end{proof}

\section{Proof by cases}

To prove $\varphi$, split the situation into finitely many exhaustive cases
and prove $\varphi$ in each --- justified by
$(p_1 \vee \cdots \vee p_n) \wedge \bigwedge_i (p_i \to q) \models q$.

\begin{theorem}
For every integer $n$, $n^{2} \geq n$.
\end{theorem}

\begin{proof}
Case $n \leq 0$: then $n^2 \geq 0 \geq n$. Case $n \geq 1$: multiplying
$n \geq 1$ by $n > 0$ gives $n^{2} \geq n$. The cases are exhaustive.
\end{proof}

\section{Existence proofs}

An existence statement $\exists x\,P(x)$ may be proved
\emph{constructively} --- exhibit a witness --- or
\emph{non-constructively} --- show a witness must exist without naming one.

\begin{example}[Constructive]
There is a positive integer expressible as a sum of two cubes in two ways:
$1729 = 1^3 + 12^3 = 9^3 + 10^3$.
\end{example}

\begin{example}[Non-constructive]\label{ex:noncon}
There exist irrationals $x, y$ with $x^{y}$ rational. Consider
$z = \sqrt2^{\sqrt2}$. If $z$ is rational, take $x = y = \sqrt2$. If $z$ is
irrational, take $x = z$, $y = \sqrt2$: then
$x^{y} = \sqrt2^{2} = 2$. Either way a witness exists --- yet we have not
determined which.
\end{example}

\section{Uniqueness proofs}

To prove $\exists!\,x\,P(x)$: prove existence, then prove uniqueness ---
assume $P(a)$ and $P(b)$ and derive $a = b$.

\begin{example}
The equation $ax = b$ ($a \neq 0$) has the unique solution $x = b/a$:
existence by substitution; uniqueness since $ax_1 = b = ax_2$ implies
$x_1 = x_2$ after dividing by $a \neq 0$.
\end{example}

\section{Disproof by counterexample}

To refute $\forall x\,P(x)$, one counterexample suffices.

\begin{example}
``Every positive integer is a sum of two squares'' is false: $3$ is not.
``$2^{2^n} + 1$ is prime for all $n$'' (Fermat) fails at $n = 5$:
$2^{32} + 1 = 641 \times 6700417$ (Euler).
\end{example}

\section{Common errors}

\begin{notebox}
Beware: proving $p \to q$ by assuming $q$ (begging the question); arguing
from examples (``it works for $n = 1,2,3$, hence always''); dividing by a
quantity that may be zero; affirming the consequent; and the hidden
assumption that objects named differently are distinct. A classic bogus
proof that $1 = 2$ divides both sides of $a^2 - ab = a^2 - b^2$ by $a - b$
when $a = b$ --- division by zero in disguise.
\end{notebox}

\section{Strategy}

Given a conjecture $p \to q$: try \emph{direct} first; if $\neg q$ unpacks
better, try \emph{contraposition}; if both directions stall or the statement
is a pure negative (``is irrational'', ``no largest \ldots''), try
\emph{contradiction}; if natural subcases appear (parity, sign, remainders),
use \emph{cases}. For $\exists$, look for an explicit witness before
resorting to non-constructive arguments. And before proving anything hard,
hunt briefly for a counterexample --- failure to find one often reveals the
proof idea. Statements quantified over all $n \in \N$ usually call for
\emph{induction}, the subject of Chapter~13.

\section{Summary}

Direct, contrapositive, contradiction and cases prove implications;
witnesses or indirect arguments prove existence; existence plus
``two solutions are equal'' proves uniqueness; one counterexample refutes a
universal claim. The two jewels --- irrationality of $\sqrt2$
(Theorem~\ref{thm:sqrt2}) and Euclid's infinitude of primes
(Theorem~\ref{thm:euclid}) --- model the contradiction template every
mathematician carries.

\section*{Exercises}
\addcontentsline{toc}{section}{Exercises}

\begin{exercise}{\easy}
Prove directly: the sum of two odd integers is even.
\end{exercise}

\begin{exercise}{\easy}
Prove directly: if $a \mid b$ and $b \mid c$ then $a \mid c$.
\end{exercise}

\begin{exercise}{\easy}
Find counterexamples: (a) ``every prime is odd''; (b) ``if $n^2 > 0$ then
$n > 0$''; (c) ``the sum of two irrationals is irrational''.
\end{exercise}

\begin{exercise}{\medium}
Prove by contraposition: if $3n + 2$ is odd, then $n$ is odd.
\end{exercise}

\begin{exercise}{\medium}
Prove by contradiction: there is no smallest positive rational number.
\end{exercise}

\begin{exercise}{\medium}
Prove by cases: for all integers $n$, $3 \mid n(n+1)(n+2)$.
\end{exercise}

\begin{exercise}{\medium}
Prove that $\sqrt{3}$ is irrational, adapting the proof of
Theorem~\ref{thm:sqrt2}. Where would the same argument break for
$\sqrt{4}$?
\end{exercise}

\begin{exercise}{\medium}
Prove there exists a unique real $x$ with $x^{3} + x - 2 = 0$.
\emph{(Existence: exhibit it. Uniqueness: factor or use monotonicity.)}
\end{exercise}

\begin{exercise}{\hard}
Prove that if $n$ is a positive integer and $n^{2}$ is divisible by $3$,
then $n$ is divisible by $3$ (use cases on $n \bmod 3$); this is the lemma
your $\sqrt3$ proof needed.
\end{exercise}

\begin{exercise}{\hard}
Show non-constructively that in any group of $2$ or more people at a party,
two people have shaken the same number of hands (handshakes within the
group). \emph{Hint: pigeonhole on the possible counts.}
\end{exercise}

\begin{exercise}{\hard}
Criticise this ``proof'' that every positive integer $n$ equals the next:
``$n = n+1$ implies $(n)^2 - (n+1)n = n^2 - (n+1)n$, both sides simplify to
$-n = -n$, true; hence $n = n + 1$.'' Identify the logical error precisely.
\end{exercise}

\section*{Solutions to the Exercises}
\addcontentsline{toc}{section}{Solutions to the Exercises}

\begin{solution}{12.1}
Let $m = 2j+1$, $n = 2k+1$. Then $m + n = 2(j + k + 1)$, an even number.
\end{solution}

\begin{solution}{12.2}
$b = aj$ and $c = bk$ for integers $j,k$; then $c = a(jk)$, so $a \mid c$.
\end{solution}

\begin{solution}{12.3}
(a) $2$. (b) $n = -1$: $n^2 = 1 > 0$ but $n < 0$. (c) $\sqrt2 + (-\sqrt2)
= 0 \in \Q$.
\end{solution}

\begin{solution}{12.4}
Contrapositive: if $n$ is even, $3n + 2$ is even. Let $n = 2k$:
$3n + 2 = 6k + 2 = 2(3k+1)$, even. Done.
\end{solution}

\begin{solution}{12.5}
Suppose $r > 0$ were the smallest positive rational. Then $r/2$ is rational,
positive, and smaller --- contradiction.
\end{solution}

\begin{solution}{12.6}
Among any three consecutive integers one is divisible by $3$: by cases on
$n \bmod 3$. If $n \equiv 0$, then $3 \mid n$; if $n \equiv 1$, then
$3 \mid n + 2$; if $n \equiv 2$, then $3 \mid n + 1$. In each case the
product is divisible by $3$.
\end{solution}

\begin{solution}{12.7}
Suppose $\sqrt3 = a/b$ in lowest terms; then $a^2 = 3b^2$, so $3 \mid a^2$,
hence $3 \mid a$ (Exercise 12.9): $a = 3c$, so $3b^2 = 9c^2$,
$b^2 = 3c^2$, and $3 \mid b$ --- contradicting lowest terms. For $\sqrt4$,
the lemma fails: $4 \mid a^2$ does \emph{not} force $4 \mid a$ (take
$a = 2$), and indeed $\sqrt4 = 2$ is rational.
\end{solution}

\begin{solution}{12.8}
Existence: $x = 1$ works: $1 + 1 - 2 = 0$. Uniqueness:
$x^3 + x - 2 = (x-1)(x^2 + x + 2)$, and $x^2 + x + 2 =
(x + \tfrac12)^2 + \tfrac74 > 0$ has no real roots; alternatively
$f(x) = x^3 + x - 2$ is strictly increasing
($f'(x) = 3x^2 + 1 > 0$), so it crosses zero exactly once.
\end{solution}

\begin{solution}{12.9}
Cases. If $n = 3k$: done, $3 \mid n$. If $n = 3k+1$:
$n^2 = 9k^2 + 6k + 1 \equiv 1 \pmod 3$, not divisible. If $n = 3k+2$:
$n^2 = 9k^2 + 12k + 4 \equiv 1 \pmod 3$, not divisible. So $3 \mid n^2$
occurs only in the first case, where $3 \mid n$.
\end{solution}

\begin{solution}{12.10}
With $n$ people, each handshake count lies in $\{0, 1, \ldots, n-1\}$ ---
$n$ values. But counts $0$ and $n-1$ cannot both occur (someone who shook
everyone's hand rules out someone who shook none). So at most $n - 1$
values are available for $n$ people: by pigeonhole two people share a
count. No specific pair is exhibited --- the proof is non-constructive.
\end{solution}

\begin{solution}{12.11}
The argument assumes $n = n + 1$ \emph{first} and derives a true statement.
Deriving a truth from an assumption proves nothing about the assumption
(that is affirming the consequent at the level of proofs); a valid proof
must derive the \emph{conclusion} from accepted premises, not the reverse.
The implication chain is also reversible only if every step is an
equivalence, which squaring-style manipulations are not.
\end{solution}
